\section{Setting: model, observations, estimators and instruments}
\label{sec:setting}

\subsection{The simulator}
\label{sec:model}

We use a differentiable Bergman-type glucose--insulin model implemented in JAX with adaptive-step Tsit5
integration (diffrax), in double precision, and a two-compartment gut. After a meal at time $0$, with
plasma glucose $G$, plasma insulin $I$, remote insulin action $X$, and stomach and intestinal glucose
$Q_1$ and $Q_2$,
\begin{align}
\dot Q_1 &= -\ke\, Q_1 + m(t), &
\dot Q_2 &= \ke\, Q_1 - \ka\, Q_2, \label{eq:gut}\\
\dot G &= -(p_1 + X)\,G + p_1 G_b + \frac{\ka\, Q_2}{V_g}, &
\dot I &= -n\,(I - I_b) + \gamma\,[G - G_b]_+, \label{eq:glucose}\\
\dot X &= -p_2 X + p_3\, \SI\,(I - I_b), \label{eq:action}
\end{align}
where $m(t)$ is a Gaussian pulse (width $2$\,min) delivering the glucose mass of the meal,
$V_g=1.6\,W$ dL for body weight $W$, and $G_b$, $I_b$ are basal values. The positive part is smoothed by a
softplus. The full engine carries further slow substrate states (glycogen, ketones, energy expenditure);
in the postprandial window the glucose dynamics are carried by the five states of
Eqs.~\eqref{eq:gut}--\eqref{eq:action}, not by three. The complete specification, with every constant and
the code of the vector field, is exported from the code and given in Appendix~\ref{app:spec}.

Each person has three free parameters, $\theta=(\SI,\ke,\ka)$, each confined to a box
(Table~\ref{tab:bounds}); everything else is fixed at population values. The \emph{primary box} is the
one fixed in the analysis plan; we also report a box half as wide and one twice as wide, scaled in
logarithms about the centre of each interval. The gut starts empty at the moment of each meal
(Section~\ref{sec:theory} explains why this convention matters).

\begin{table}[t]
\centering\small
\begin{tabular}{lcccl}
\toprule
parameter & symbol & population value & primary box (min$^{-1}$ for rates) & role \\
\midrule
insulin sensitivity & $\SI$ & $1$ & $[0.3,\,1.6]$ & scales insulin action \\
gastric emptying & $\ke$ & $0.026$ & $[0.015,\,0.04]$ & stomach $\to$ intestine \\
carbohydrate absorption & $\ka$ & $0.022$ & $[0.012,\,0.032]$ & intestine $\to$ plasma \\
\bottomrule
\end{tabular}
\caption{The three personalised parameters and their primary parameter box. All other parameters are
fixed.}
\label{tab:bounds}
\end{table}

\subsection{Observables}
\label{sec:observables}

A \emph{meal record} is a window of CGM samples around one meal, with the logged carbohydrate, fat and fibre.
All observables are computed from the baseline-subtracted excursion $y(t)=[G(t)-\bar G_{\mathrm{pre}}]_+$,
where $\bar G_{\mathrm{pre}}$ is the mean of the samples in the $30$ minutes before the meal.
\begin{itemize}
\item \emph{iAUC}: $\int y\,dt$ over the post-meal window. One number per meal. This is the target of the
  primary analysis.
\item \emph{Centroid}: $\int t\,y\,dt \,/ \int y\,dt$, the time centre of mass of the excursion. One number
  per meal, carrying timing and almost no magnitude.
\item \emph{Trace}: the baseline-subtracted curve itself, at the sampling resolution of the cohort.
\end{itemize}
The observables form a \emph{ladder}: iAUC, iAUC with centroid, full trace. Each has a smooth form (used
for gradients; softplus and softmax of sharpness $10$) and a hard form (used for scoring); the two agree to
within one percent of the hard value on real excursions.

\subsection{Likelihood, estimation and the identifiability criterion}
\label{sec:estimation}

For every identifiability analysis the estimate is the \emph{unregularized maximum-likelihood} estimate,
never the penalized prediction fit, because a penalty toward the population values would narrow every
interval. Residuals are Gaussian with a per-subject, per-observable variance estimated at a pilot fit and
then held fixed; stacked observables are weighted by their inverse residual variance; and the trace is
whitened for first-order autocorrelation of the sensor residuals (Prais--Winsten). Freezing the noise
parameters matters: re-estimating them inside the optimization would let the fit lower its loss by
declaring its errors noise. With the negative log-likelihood $\NLL(\theta)$, a $95\%$ profile interval
for $\theta_j$ is the set where
$\min_{\theta_{-j}}\NLL(\theta_j,\theta_{-j})-\min_\theta \NLL(\theta) \le \delta$, with
$\delta=\resDeltaThreshold$ (half the $0.95$ quantile of $\chi^2_1$).

We classify each profile as \emph{bounded} (the profile rises above $\delta$ on both sides inside the box),
\emph{one-sided} (on one side only) or \emph{flat} (never reaches $\delta$ inside the box). Only a bounded
interval is called identified; a one-sided interval is constrained by the box on the other side and
therefore by the prior range, which is why every result is also reported at three box widths. For the two
rates only, the grid is additionally extended to four times the upper bound, as a diagnostic that is kept
out of the classification.

A parameter's interval can be bounded only if the estimate is not on a bound. Subject~$i$ is
\emph{interior for parameter $j$} when the estimate lies more than $1\%$ of the box width from both bounds,
and fractions of bounded intervals are reported both over all subjects and over subjects interior for the
parameter. (Defining interior as ``no parameter on a bound'' gives an empty set whenever one parameter is
unidentified, a defect of the original plan documented in Appendix~\ref{app:register}.)

\subsection{Instruments}
\label{sec:instruments}

\paragraph{Fisher information.}
We compute $J=\partial(\text{weighted predictions})/\partial\log\theta$ by reverse-mode automatic
differentiation, $F=J^\top J$, its eigenvalues and eigenvectors, and the Schur complement of the timing
block $\{\ke,\ka\}$ after profiling out $\SI$, whose condition number measures how well the data
constrain the timing parameters jointly. The \emph{exchange direction} is
$(1/\ka,\,-1/\ke)$ in $(\log\ke,\log\ka)$: the direction along which $\tau_1=1/\ke+1/\ka$ is constant to
first order. We report the absolute cosine between the weakest eigenvector of the timing block and it.

\paragraph{Generic rank.}
The rank of the stacked sensitivity matrix over many random parameter draws, at a relative tolerance of
$10^{-6}$, distinguishes a parameter set that is locally identifiable but ill-conditioned from one that is
rank deficient.

\paragraph{Structural identifiability.}
Symbolic analysis with StructuralIdentifiability.jl, under several formulations that differ in what is
assumed known about the gut at the start of a meal (Section~\ref{sec:structural}).

\paragraph{Replica and synthetic study.}
In the \emph{replica}, each subject's regularized prediction-fit estimate (moved at least $5\%$ of the box
width inside the box) is taken as the truth; each real meal is simulated at it; autoregressive CGM noise
with the subject's own variance and lag-one correlation (from the residual of the real fit) is added; and
the carbohydrate given to the fit is the true value multiplied by lognormal error of unit mean and
coefficient of variation $0.25$. The replica answers: if the model were true, with this noise and this
logging error, how often would the data bound a parameter? In the \emph{synthetic study}, the true
parameters are instead drawn independently and uniformly inside the box, so that recovery can be read
against the location of the truth. In both, the truth is known and the coverage of bounded intervals can be
checked. Random streams are seeded from a hash of the subject, the replica seed and the stream name, so
switching one noise source off does not change the other's draws.

\paragraph{Optimizer check.}
Every estimate comes from fixed-step Adam in log coordinates. For the rate-space fits we compare it with
L-BFGS-B from five random starting points (the gap in negative log-likelihood, to be read against
$\delta$); for the coordinate fits of Section~\ref{sec:more} the L-BFGS-polished estimate replaces it.

\subsection{Cohorts}
\label{sec:cohorts}

Table~\ref{tab:cohorts} lists the cohorts, which differ in ways that matter: the sensor interval, the
observation window, and how carbohydrate was obtained. We use subjects with at least ten meals (five for the
cohort with standardized meals). In the Shanghai data the sensor samples every $15$ minutes and the stored
$5$-minute grid was filled by interpolation, so every trace objective reads only the real samples.

\begin{table}[t]
\centering\small
\begin{tabular}{lrrcccr}
\toprule
cohort & subjects & meals & meals per subject & sensor (min) & window (min) & median carbs (g) \\
\midrule
CGMacros & 45 & 1640 & 36 (17--64) & 5 & 180 & 56 \\
Shanghai & 87 & 2608 & 30 (10--56) & 15 & 180 & 33 \\
Standardized meals & 15 & 82 & 5 (5--6) & 5 & 145 & 35 \\
\bottomrule
\end{tabular}
\caption{The cohorts after the selection rule (at least ten meals per subject, five for the standardized-meal cohort). Meals per subject: median (minimum--maximum). The Shanghai sensor samples every fifteen minutes; its stored five-minute grid is interpolated and trace objectives read the real samples only. Carbohydrate sources are described in the text.}
\label{tab:cohorts}
\end{table}


Carbohydrate comes from a photographed, nutritionist-coded log in CGMacros; from a free-text dietary record
in the Shanghai cohort, from which we estimate carbohydrate by keyword from the food name and weight (an
estimate with its own error, which is one reason for the carbohydrate-error experiments of
Section~\ref{sec:si}); and from the study design in the standardized-meal cohort, whose meals are known by
construction. The window is the observation window of the iAUC, which is shorter in the last cohort.

\subsection{Prediction protocol}
\label{sec:cv}

Held-out prediction uses $5$-fold cross-validation within subject, repeated $5$ times with SHA-256-seeded
fold assignments identical for every method, so that no method benefits from a favourable split; one unit is a (subject, repeat), every meal is held out once, a
subject's score is its mean over repeats, and every interval resamples \emph{subjects}. Cells compared:
gradient fitting with three parameters or with $\SI$ only, exhaustive grid search of the same loss (a
$20\times8\times8$ grid, or a $41$-point grid), the same under the full-trace objective, a random forest
and a simulation-based-inference baseline trained on simulated data, and the personal mean, population
default and persistence baselines. Equivalence of two methods is declared when the $90\%$ bootstrap interval
of their paired difference lies inside $(-150,+150)$ mg\,dL$^{-1}$\,min (margins $90$ and $300$ also
reported); a difference is not called equivalence unless that test was run.

\subsection{Statistics and the analysis plan}
\label{sec:plan}

Fractions of subjects carry Wilson intervals. Paired differences carry subject-level BCa bootstrap
intervals ($\resBootstrapResamples$ resamples), Wilcoxon tests with Holm correction within families, and
two one-sided tests for equivalence. The analysis plan, with thresholds, was fixed in version control before
the analyses were run (hypotheses H1--H10); amendments were committed before the analyses they concern or,
where stated, after results had been seen. Hypotheses H11--H12 were declared before their analyses;
\textbf{H13--H22 were declared after earlier results were seen} and are marked post hoc throughout. Every
hypothesis, its rule, its status and the deviations are listed in Appendix~\ref{app:register}; the
conclusions below do not depend on any threshold having been changed, because none was.
